% META: {"id": "TCOMPL-ATT-EX-002", "chapitre": "TCOMPL-TEMPS-ATTENTE", "type_objet": "exercice", "capacites": ["TCOMPL-ATT-C3"], "capacites_codes": ["C3"], "methodes": [], "parcours": 2, "competences": ["raisonner"], "duree_min": 15, "mode_creation": "ex_nihilo", "sources_inspiration": [], "fichier_tex": "chapitres/TCOMPL-TEMPS-ATTENTE/exercices/TCOMPL-ATT-EX-002.tex", "corrige_tex": "chapitres/TCOMPL-TEMPS-ATTENTE/corriges/TCOMPL-ATT-CO-002.tex", "status": "generated"}
% BEGIN-VERIFY
% from sympy import Rational
% p = Rational(1, 4)
% def PX_sup(n, p):
%     return (1-p)**n
% n, k = 4, 6
% cond = PX_sup(n+k, p) / PX_sup(n, p)
% sans_memoire = PX_sup(k, p)
% assert cond == sans_memoire
% assert sans_memoire == (Rational(3,4))**6
% END-VERIFY
\begin{exercice}{TCOMPL-ATT-EX-002}{2}{15}
Un standard téléphonique répond au premier appel avec une probabilité de $25\%$ à chaque tentative. On sait que les $4$ premières tentatives ont échoué.
\begin{enumerate}
  \item En utilisant l'absence de mémoire, exprimer la probabilité d'attendre encore $6$ tentatives supplémentaires avant le succès, sans refaire de calcul complexe.
  \item Calculer cette probabilité.
\end{enumerate}
\end{exercice}
